% !TEX root = ../algebra_02.tex

\section{Линейные операторы. Инвариантные подпространства}

\begin{Definition}{Линейный оператор}{}
	Линейным оператором на $V$ называется линейное отображение из пространства $V$ в себя. Множество всех таких операторов обозначается $\End V:=\Hom(V,V)$.
\end{Definition}

Важную роль в изучении структуры линейного оператора играют подпространства, которые оператор переводит в себя.

\begin{Definition}{Инвариантное подпространство}{}
	Пусть $\alpha\in \End V$ и $W\le V$. Подпространство $W$ называется инвариантным относительно $\alpha$, если
	\[
	\alpha W\subset W
	\].
\end{Definition}

Наличие инвариантного подпространства позволяет упростить матрицу оператора при правильном выборе базиса.

\begin{Theorem}{Блочный вид при инвариантном подпространстве}{}
	Пусть $W\le V$, $\dim W=m$, и $E=(e_1,\ldots,e_n)$ --- базис $V$ такой, что первые векторы $e_1,\ldots,e_m$ образуют базис $W$. Тогда $W$ инвариантно относительно $\alpha$ тогда и только тогда, когда матрица оператора имеет блочный вид:
	\[
	[\alpha]_E=
	\begin{pmatrix}
		B & *\\
		0 & C
	\end{pmatrix},
	\]
	где $B\in M_m(K)$.
\end{Theorem}

\begin{proof}
	Если $W$ инвариантно, то для любого $j=1,\ldots,m$ имеем $\alpha e_j\in W$. Значит, в координатах по базису $E$ у вектора $\alpha e_j$ последние $n-m$ координат равны нулю. Именно поэтому в первых $m$ столбцах матрицы снизу стоит нулевой блок.

	Обратно, если матрица имеет такой блочный вид, то для каждого $j=1,\ldots,m$ образ вектора $\alpha e_j$ выражается только через $e_1,\dots,e_m$, то есть лежит в $\langle e_1,\ldots,e_m\rangle=W$. Следовательно, $\alpha W\subset W$.
\end{proof}

\begin{Corollary}{Инвариантность прямых слагаемых}{}
	Пусть $V=W_1\oplus W_2$,
	$e_1,\ldots,e_m$ --- базис $W_1$, а $e_{m+1},\ldots,e_n$ --- базис $W_2$.
	Для $E=(e_1,\ldots,e_n)$ следующие условия эквивалентны:
	
	\begin{enumerate}
		\item $W_1$ и $W_2$ инвариантны относительно $\alpha$;
		\item матрица оператора имеет блочно-диагональный вид
		\[
		[\alpha]_E=
		\begin{pmatrix}
			B & 0\\
			0 & C
		\end{pmatrix}.
		\]
	\end{enumerate}
\end{Corollary}
